\begin{explanationbox}

	\textbf{\textcolor{CExplanation}{一句话直觉}}

	投影数量是向量 $\vec a$ 在 $\vec b$ 方向上的有向影子长度：它是标量，可以为正、为负或为零。

	\medskip
	\textbf{\textcolor{CExplanation}{核心拆解}}
	\begin{description}[style=nextline, leftmargin=2.8em, labelsep=0.5em, itemsep=0.45em]
		\item[\textbf{要点一（数量公式）：}] 当 $\vec b\neq\vec 0$ 时，$\vec a$ 在 $\vec b$ 方向上的投影数量为
		      \[
			      \frac{\vec a\cdot\vec b}{\lvert\vec b\rvert}.
		      \]
		      若 $\vec a\neq\vec 0$，且 $\theta$ 是 $\vec a$ 与 $\vec b$ 的夹角，则也可写成
		      \[
			      \lvert\vec a\rvert\cos\theta .
		      \]

		\item[\textbf{要点二（符号判定）：}] 当 $\vec a\neq\vec 0,\ \vec b\neq\vec 0$ 时：
		      \[
			      0 \le \theta < \frac{\pi}{2}
			      \quad\Rightarrow\quad
			      \text{投影数量为正};
		      \]
		      \[
			      \theta = \frac{\pi}{2}
			      \quad\Rightarrow\quad
			      \text{投影数量为零};
		      \]
		      \[
			      \frac{\pi}{2} < \theta \le \pi
			      \quad\Rightarrow\quad
			      \text{投影数量为负}.
		      \]
		      特别注意：若 $\vec a=\vec 0$，投影数量为 $0$，但此时不能用“夹角为直角”来解释。

		\item[\textbf{要点三（点积关系）：}] 由点积定义
		      \[
			      \vec a\cdot\vec b=\lvert\vec a\rvert\lvert\vec b\rvert\cos\theta,
		      \]
		      因此在 $\vec b\neq\vec 0$ 时，
		      \[
			      \lvert\vec a\rvert\cos\theta
			      =
			      \frac{\vec a\cdot\vec b}{\lvert\vec b\rvert}.
		      \]
	\end{description}

	\medskip
	\textbf{\textcolor{CExplanation}{几何本质（有向投影长）}}

	过 $\vec a$ 的终点向 $\vec b$ 所在直线作垂线，垂足与原点的有向距离就是投影数量。其绝对值等于投影线段长，符号对应投影方向与 $\vec b$ 相同或相反。

	当 $\vec a=\vec 0$ 时，终点就是原点，投影数量为 $0$；但零向量与 $\vec b$ 的夹角不需要讨论。

	\medskip
	\textbf{\textcolor{CExplanation}{代数意义（点积归一化）}}

	点积 $\vec a\cdot\vec b$ 可看作“$\vec a$ 在 $\vec b$ 方向上的投影数量”再乘以 $\lvert\vec b\rvert$。因此
	\[
		\frac{\vec a\cdot\vec b}{\lvert\vec b\rvert}
	\]
	相当于把目标方向 $\vec b$ 的长度影响去掉，只保留 $\vec a$ 在 $\vec b$ 单位方向上的有向分量。

	\medskip
	\textbf{\textcolor{CExplanation}{考点价值（两种考法）}}
	\begin{itemize}[leftmargin=*, itemsep=0.5em, topsep=0.4em]
		\item \textbf{考法一（直接求投影）：} 已知向量坐标或模长及夹角，直接代入
		      \[
			      \frac{\vec a\cdot\vec b}{\lvert\vec b\rvert}
		      \]
		      计算。使用前必须确认 $\vec b\neq\vec 0$。

		\item \textbf{考法二（判夹角范围）：} 当 $\vec a\neq\vec 0,\ \vec b\neq\vec 0$ 时，可通过投影数量的正负反推 $\cos\theta$ 的符号，从而判断夹角是锐角、直角还是钝角。
	\end{itemize}

	\medskip
	\textbf{\textcolor{CExplanation}{顿悟点}}

	投影数量可以看成点积的“单位方向版本”：点积 $\vec a\cdot\vec b$ 同时受到 $\vec a$ 的长度、$\vec b$ 的长度以及夹角影响；而投影数量
	\[
		\frac{\vec a\cdot\vec b}{\lvert\vec b\rvert}
	\]
	相当于把目标方向 $\vec b$ 单位化后，衡量 $\vec a$ 在该单位方向上的有向分量。

	\medskip
	\textbf{\textcolor{CExplanation}{使用场景}}
	\begin{itemize}[leftmargin=*, itemsep=0.5em, topsep=0.4em]
		\item \textbf{场景一（物理做功）：} 恒力 $\vec F$ 在位移 $\vec s$ 方向上的投影数量等于 $\lvert\vec F\rvert\cos\theta$，功
		      \[
			      W=\vec F\cdot\vec s=(\text{力在位移方向上的投影数量})\times \lvert\vec s\rvert .
		      \]

		\item \textbf{场景二（向量分解）：} 将 $\vec a$ 分解为平行于 $\vec b$ 的分量
		      \[
			      \vec a_{\parallel}
			      =
			      \frac{\vec a\cdot\vec b}{\lvert\vec b\rvert^2}\vec b
		      \]
		      和垂直于 $\vec b$ 的分量。其中投影数量
		      \[
			      \frac{\vec a\cdot\vec b}{\lvert\vec b\rvert}
		      \]
		      是平行分量在 $\vec b$ 方向上的有向长度；而平行分量的实际长度为
		      \[
			      \left\lvert \vec a_{\parallel}\right\rvert
			      =
			      \frac{\lvert\vec a\cdot\vec b\rvert}{\lvert\vec b\rvert}.
		      \]
	\end{itemize}
\end{explanationbox}